%% =============================================================================
%% SUPPLEMENT S1 -- the pre-registered negative-result section, as drafted
%%
%% This is the text of "If the distinction does not separate" exactly as it
%% stood in build/main.tex BEFORE the full-sample analysis was run. It is
%% anti-tuning commitment 4 of the pre-registration (https://osf.io/ezjh2/,
%% DOI 10.17605/OSF.IO/EZJH2): its purpose was to fix, in advance and in the
%% authors' own words, what the paper would say if H1 failed, so that a null
%% result had a finished home and there was no incentive to respecify the
%% model until it passed.
%%
%% H1 was supported, so the section did not become operative. Following the
%% instruction written into the section itself before the result was known, it
%% is reproduced here verbatim and replaced in the manuscript by a short
%% subsection recording that it was prepared and not needed. It is reproduced
%% rather than deleted because a commitment that vanishes when it is not
%% needed was never a commitment.
%%
%% NOT EDITED after the verdict. The only change is this header.
%% =============================================================================

\section{If the distinction does not separate}\label{sec:negative}

%% ==========================================================================
%% THIS SECTION IS WRITTEN BEFORE THE ANALYSIS IS RUN.
%%
%% It is anti-tuning commitment 4 in the pre-registration
%% (https://osf.io/ezjh2/, DOI 10.17605/OSF.IO/EZJH2). Its purpose is to fix,
%% in advance and in the authors' own words, what the paper says if H1 fails
%% -- so that a null result has a finished home to go to and there is no
%% incentive to respecify the model until it passes.
%%
%% DO NOT DELETE THIS SECTION IF H1 IS SUPPORTED. Convert it into a short
%% paragraph reporting that the pre-registered null framing was prepared and
%% did not become necessary, and move this text to the supplement. Deleting
%% it silently would defeat its purpose.
%% ==========================================================================

Hypothesis H1 holds that individualism--collectivism and altruism--egoism are
empirically separable at population level. It is pre-registered with a
conjunctive decision rule: a two-factor model must beat a one-factor model by
$\Delta\mathrm{CFI} \geq .010$, the model fixing $\varphi(\mathrm{D1},
\mathrm{D2}) = 1$ must be rejected at $\alpha = .01$, and the estimated
$\varphi(\mathrm{D1}, \mathrm{D2})$ must not exceed $.85$. If any of the three
fails, H1 is unsupported, and this paper reports it as unsupported. We set out
here, before running the analysis, what that finding would and would not mean,
so that the interpretation is not constructed after the fact.

\subsection{What a null on H1 would establish}

It would establish that the two axes Popper distinguishes do not separate
\emph{in these indicators, in these populations}. That is a substantive result
about European publics as measured by the EVS/WVS and ESS item pools, and it is
worth reporting for three reasons.

First, the distinction is load-bearing in political theory. Popper's argument
in \emph{The Open Society and Its Enemies} \citep{popper_ose} turns on the claim that the
attractive cell -- individualist and altruistic -- is genuinely distinct from
atomism on one side and from collectivist altruism on the other. A century of
commentary has treated that four-cell table as though its axes were obviously
independent. Nobody has tested it. A negative answer constrains how the
argument may be used empirically, which is exactly the kind of finding a
measurement paper exists to produce.

Second, a null would bear directly on interpretive practice. Comparative
researchers routinely import a single individualism score as a covariate and
read substantive meaning into it. If D1 and D2 are not separable in the data,
that practice is not thereby vindicated: the finding would be that the
available instruments cannot distinguish the cells, which is a reason for
caution about what any individualism score means, not a licence to keep using
one.

Third, H1 is only one of the paper's hypotheses, and not the primary one. H2 --
that institutional respect is not reducible to individualism -- is the headline
claim, and it is logically independent of H1. A paper reporting a null on H1
and a positive on H2 is not a failed paper; it is a paper reporting that one of
Popper's two measurement-relevant distinctions survives population-level
testing and the other does not.

\subsection{What a null on H1 would not establish}

It would not establish that Popper was wrong. The inference from ``these items
do not separate'' to ``the distinction is unreal'' requires the premise that
the items adequately operationalise both axes, and we have grounds -- stated
before seeing the data -- for doubting that premise on the D2 side
specifically.

D2 is the thinner dimension in every instrument we examined. In the joint
EVS/WVS core it is carried by one directly on-target item (a child-quality
item naming unselfishness), a tolerance item, a block of justifiability items
about defection from collective goods, and two behavioural
organisation-membership items. Several of these are binary. None was designed
to measure direction of concern; they were designed for other purposes and are
being repurposed here. D1, by contrast, is measured by long, well-developed
batteries built over decades precisely to capture autonomy and choice.

An asymmetry of that size has a predictable consequence. A weakly measured
factor with few indicators and low reliability will correlate highly with a
strongly measured one for reasons that have nothing to do with the constructs:
attenuation runs the other way, but common method variance, shared response
format and the sheer instability of a short factor all inflate
$\varphi$. If H1 fails, the most likely explanation is not that Popper's
distinction is empty but that no existing cross-national survey was built to
test it. We will say so, and we will say it having committed to saying it in
advance.

A null would also not license quietly dropping D2 and reporting a
two-dimensional model as though that had been the plan. The pre-registration
fixes the item table and the model sequence; a failed H1 is reported against
the registered specification, not against a revised one.

\subsection{How we will tell the two explanations apart}

The distinction between ``the axes do not separate'' and ``our D2 does not
measure anything stable'' is partly testable, and we pre-commit to the
diagnostics rather than choosing them afterwards.

If D2's indicators show adequate internal structure -- acceptable loadings,
a reliability estimate appropriate to the number and format of the indicators,
and a factor that behaves consistently across the primary and replication
samples -- and $\varphi(\mathrm{D1}, \mathrm{D2})$ is nevertheless near unity,
then the null is about the constructs and we will report it that way.

If instead D2's loadings are weak or unstable, or the factor fails to replicate
between the EVS/WVS and ESS samples, then the null is uninformative about
Popper and we will report it as a measurement failure and as a specification
for what a purpose-built instrument would need to contain. That specification
is itself a contribution: it says what a survey would have to ask to test a
distinction the field has been assuming for seventy years.

The behavioural indicators matter here. Organisation-membership items share no
response format with the attitudinal batteries, so if the attitudinal and
behavioural indicators of D2 cohere, common method variance is a less plausible
account of any collapse; if they do not cohere, the factor is not yet a factor.

\subsection{The commitment}

We will not respecify the model to rescue H1. We will not drop indicators to
improve fit, add cross-loadings suggested by modification indices, or exchange
the pre-registered decision rule for a more forgiving one. Items may leave the
model only for the two reasons fixed in the pre-registration -- a standardised
loading below $.30$ in the configural solution, or a Heywood case -- and every
removal is logged with its reason in the supplement.

If the primary and replication samples disagree, we report a split decision as
a split decision rather than resolving it in favour of the positive result.

\subsection{The same commitment for the synthetic panel}

The panel component carries a parallel risk, and the parallel commitment. Its
evidentially load-bearing tier is a forecast of this paper's own analysis,
elicited and time-stamped before that analysis was run. If the panel forecasts
poorly -- if its Brier skill against the outcome is at or below the naive
baseline -- that is the result, and it is reported as the result. A synthetic
panel that cannot forecast a measurement model it was given the full
specification for is evidence about the limits of synthetic panels, which is
precisely what this part of the paper is for.

Nor will a good Brier score be read as validation of the panel in general. The
forecast covers one analysis, in one literature, with a particular item pool.
The panel remains an instrument with a documented error model. It is not a
substitute for expert consultation, and no result reported here should be taken
to suggest otherwise.
